\section{Where the Rules Came From}
The argument requires three things of history, and this section establishes each from the paper's own evidence: the rules were written by nameable coalitions; they have barely moved since; and they leaked from birth. Table~\ref{tab:genealogy} summarises the genealogy for five focal states.

Written by nameable coalitions. These rules do not have anonymous origins. Florida's constitution once stated who the consenting public was in so many words: freeholders, property owners, the only electors permitted to vote on bonds [V: article and date]. California's two-thirds requirement and debt limit entered together in the 1879 constitution, the work of a convention dominated by taxpayer and anti-corporate coalitions; Kentucky's tax-rate limits and voter-assent requirement are Sections 157 and 158 of the 1891 constitution, a taxpayer convention's settlement. The wave itself was a response to identifiable events: eleven states replaced their constitutions between 1842 and 1852 after the state debt crisis, rewriting debt procedure and corporate chartering together (Wallis 2005), and the municipal defaults after 1873 pushed the same settlement one tier down (Dove 2014; Hillhouse 1936). The referendum itself entered this history as a promoter's tool before it became a taxpayer's shield: mid-century county votes were the instrument by which communities granted railroad aid, and when the aid soured, the federal courts forced payment over state objection (Gelpcke v. Dubuque 1863) and the constitutional response converted the enabling device into a barrier (Cole 2023 [V]). In several Southern states the authoring conventions' disenfranchising purposes were explicit, and the fiscal articles carried them [V: Kousser; Alabama 1901]. The coalitions, in short, are in the record: taxpayers, creditors made cautious, and property holders whose names some texts preserved.

Frozen by comparison. The claim that the rules have not moved is measured, not asserted. Between the only two national compilations of local debt restrictions ever produced, in 1961 and 1986, 117 comparable entries yield twelve substantive legal changes; the rest of the flagged differences are definitional (source-diff table, replication files). Over the same decades the state constitutions surrounding these articles were amended freely (Dinan 2006), so the fiscal articles are frozen relative to their own documents, not merely old. The one modern movement of a threshold proves the point by its shape: Proposition 39 (2000) cut California's school bond requirement from two thirds to fifty-five per cent, a purpose-specific amendment won by the single sector the exits never reached, and the state's own analyst later quantified what the old bar had been doing: at the two-thirds requirement, a further third of school bond measures, nine billion dollars' worth between 1986 and 2000, had majority support and failed; after the change, approval rates rose from 55 to 80 per cent (California Legislative Analyst's Office 2009). The bar was binding, and everyone who could measure it knew.

Leaking from birth. The exits are not modern corrosion; they are nearly as old as the gates. Within decades of the 1879 settlement, California courts had constructed the special fund doctrine, under which obligations payable solely from project revenues are not ``debt'' within the constitutional meaning, and the lease exceptions followed [V: Offner-Dean line]. Florida's first judicial exit opened within a few years of the freeholder clause itself [V: State v. City of Miami]. By 1958 a signed Yale Law Journal article could describe public building authorities, matter-of-factly, as ``the costly subversion of state constitutions'' (Morris 1958). Legislatures wrote the leak into the statutes directly: Iowa's code authorises ``essential corporate purpose'' borrowing with no election while requiring sixty per cent approval for ``general corporate purposes,'' the scope line drawn in purpose language inside a single chapter (Iowa Code \S\S384.24 to 384.26). Kentucky ran the whole sequence in one state: the 1891 gates, a century of doctrinal exits through holding companies and revenue devices, formal amendment ratifying what practice had built, and a measurable endpoint, for in the modern corpus Kentucky's voters authorise essentially none of the state's local borrowing, the lowest value on the consent map (Figure~\ref{fig:map}). The oldest strong gate and the emptiest modern ballot are the same case. That is this section's thesis in one state, and it is why the persistence mechanism of Section 2.4 is historical before it is contemporary: the rules never accumulated the reform pressure that would have moved them, because from nearly the beginning, everyone who could leave, left.

% owner-supplied genealogy (Section 3); [V] cells to resolve against the origins file
\begin{table}[!htbp]\centering\footnotesize
\caption{The genealogy of the focal states' rules}\label{tab:genealogy}
\begin{tabular}{@{}>{\raggedright\arraybackslash}p{0.09\linewidth}>{\raggedright\arraybackslash}p{0.16\linewidth}>{\raggedright\arraybackslash}p{0.13\linewidth}>{\raggedright\arraybackslash}p{0.17\linewidth}>{\raggedright\arraybackslash}p{0.15\linewidth}>{\raggedright\arraybackslash}p{0.18\linewidth}@{}}
\toprule
State & Rule (current) & Origin & Authoring coalition & Veto-holder as written & First major exit \\
\midrule
California & 66.7\% GO; 55\% school (2000) & Const.\ 1879, art.\ [V] & Taxpayer / anti-corporate convention & Two-thirds of electors & Special fund doctrine [V: date] \\ \addlinespace
Kentucky & Voter assent above \S 157 limits & Const.\ 1891, \S\S 157--158 & Taxpayer convention & Two-thirds of voters [V: confirm share] & Holding-company / revenue devices [V] \\ \addlinespace
Texas & Simple majority (GO, by class) & Const.\ 1876, art.\ XI [V] & Post-Reconstruction retrenchment convention [V] & Majority of qualified voters & Statutory districts; landowner-franchise districts (upheld, \emph{Ball v.\ James} 1981) \\ \addlinespace
Wisconsin & Majority; ch.\ 67 referenda & Stat.\ ch.\ 67 [V: origin date] & [V] & Majority of electors & Petition and exemption classes [V] \\ \addlinespace
Louisiana & Majority; bond commission & Const.\ [V: date] & [V] & Majority of electors voting & State commission channel [V] \\
\bottomrule
\end{tabular}
\begin{minipage}{0.96\linewidth}\vspace{2pt}\scriptsize
\emph{Notes:} Rows complete to the paper's archived statutory texts; [V] cells resolve against the origins file before circulation. The table's point survives its pending cells: each rule has a date, a coalition, and a named public.
\end{minipage}
\end{table}

One correction to prior measurement closes the section. The standard cross-state codings treat these rules as a binary, referendum or not, coded from statute at the state level. The exercise above shows why that fails: a ceiling with a vote to exceed it is not a mandatory referendum; a petition right is not an election requirement; a town meeting is a voter gate wearing a council's clothes; and Iowa's purpose-split statute is both regimes at once. The historical coding that treated permission as equivalent to obligation (Dove 2014, addressed in the appendix) conflated the two things this section has distinguished, which is the difference the paper's observed-authorisation measure exists to respect.

